\section{策略融合与调度}

\subsection{组合提示、搜索与 trace}
Chen 强调上述策略不是孤立：可以先用 Tree-of-Thought 拓展 branch，再通过 self-consistency 投票稳定 final answer，最后让 trace feedback 迭代验证。换句话说，pipeline 形如 prompt → search → iteration，每一环都有对应的指标与 budget。

\subsection{预算调度与指标}
实际部署需对 token、branch、iteration 三重预算做调度。复杂任务可在前期分配更多 token，在 search 阶段用 verifier 控制 branch 数量，在 iteration 阶段限制 feedback rounds。指标上要同时记录 accuracy、token count、latency。

\begin{knowledgebox}{预算调度原则}
1. 自动伸缩：用 task difficulty signal 调整 token/branch depth；2. 优先级队列：借助 verifier 给每个 branch 打分，控制展开顺序；3. 观测表征：记录每个模块（提示、采样、trace）在 deployment 中的 F1、token count 与 latency，以便随时重新权衡资源。
\end{knowledgebox}

\begin{importantbox}{策略融合的实践提示}
不要把 prompt、search、trace 视为独立系统，而是把它们当成一个预算链条：当 search depth 增长时，自动扣减 token budget 或限制 iteration；当 iteration 深度增加时，把 search branching 退回到安全值。这样才能在部署中避免“所有模块独立延长，预算指数级膨胀”的现象。
\end{importantbox}

\subsection{本章小结}
策略融合意味着把提示、搜索、trace 看作一条闭环 pipeline。通过预算调度、优先级队列与指标回路，可以在 inference-time compute 下保持高准确率与可解释性。

